% ------------------------------------------------------------------------
%
\begin{frame}{Proving associativity of stack concatenation}

The definition of the concatenation of two stacks was:
\begin{equation*}
\fun{cat}(        \el,t) \xrightarrow{\smash{\alpha}} t;\qquad
\fun{cat}(\cons{x}{s},t) \xrightarrow{\smash{\beta}}
\cons{x}{\fun{cat}(s,t)}.
\end{equation*}
Let us prove the associativity of `\fun{cat}', which we express
formally as
\begin{equation*}
  \pred{CatAssoc}{s,t,u} \colon
\fun{cat}(s,\fun{cat}(t,u)) \equiv
\fun{cat}(\fun{cat}(s,t),u)
\end{equation*}
for all stack values \(s\), \(t\) and~\(u\).

\end{frame}

% ------------------------------------------------------------------------
%
\begin{frame}{Proving associativity of stack concatenation}

  We apply the well\hyp{}founded induction principle to the structure
  of~\(s\), so we must establish
  \begin{itemize}

    \item the basis \(\forall t,u \in S.\pred{CatAssoc}{\el,t,u}\), and the

    \item step \(\forall s,t,u \in S.\pred{CatAssoc}{s,t,u}
      \Rightarrow \forall x \in T.\pred{CatAssoc}{\cons{x}{s},t,u}\).

  \end{itemize}

  \bigskip

  Underlining the call being rewritten, the base case is:
  \begin{equation*}
    \ufun{cat}(\el,\fun{cat}(t,u))
    \xrightarrow{\smash{\alpha}} \fun{cat}(t,u)
    \xleftarrow{\smash{\alpha}} \fun{cat}(\ufun{cat}(\el,t),u).
  \end{equation*}

  Let us assume now \(\pred{CatAssoc}{s,t,u}\), called the
  \textbf{induction hypothesis}, and let us prove
  \(\pred{CatAssoc}{\cons{x}{s},t,u}\), for any term~\(x\).

\end{frame}

% ------------------------------------------------------------------------
%
\begin{frame}{Proving associativity of stack concatenation}

  The goal here is to use the rewrite system as an abstract machine to
  prove a property, with the timely help of the induction principle.

  \bigskip

  Precisely, we want to rewrite each side of the equivalence we wish
  to prove until we either find the same term (equality), or we use
  the induction hypothesis (equivalence).

  \bigskip

  Since we are going to work with equivalences, we should lift the
  call\hyp{}by\hyp{}value reduction strategy here, as it does not
  matter when arguments are evaluated, as long as their evaluation
  terminates and is confluent: we always assume these properties on
  the system hold when proving properties of the functions it defines.

\end{frame}

% ------------------------------------------------------------------------
%
\begin{frame}{Proving associativity of stack concatenation}

  We may start with the left\hyp{}hand side of
  \(\pred{CatAssoc}{\cons{x}{s},t,u}\), that is
  \(\fun{cat}(\cons{x}{s},\fun{cat}(t,u))\).

  \bigskip

  We notice that this is \textbf{not} a value, as the call contains a
  reducible call, \(\fun{cat}(t,u)\).

  \bigskip

  When using functions to compute, we match a call containing only
  values against the patterns of the rules.

  \bigskip

  Here, we would like to match a pattern against a pattern. This is
  more general than matching, because values are patterns, and is
  called \textbf{unification}. Given two patterns, we want to find
  whether we can find assignments to the variables of these so that,
  when they are substituted by their denotation, the two patterns
  become equal.

\end{frame}

% ------------------------------------------------------------------------
%
\begin{frame}{Proving associativity of stack concatenation}

 Here, we would like to unify
 \(\fun{cat}(\cons{x}{s},\fun{cat}(t,u))\) (called the \textbf{goal})
 and \(\fun{cat}(\cons{x}{s},t)\) (the pattern).

 \bigskip

The problem is that we have variables that are the same in the goal
and the pattern (\(x\), \(s\)~and~\(t\)).

\bigskip

The case of \(\cons{x}{s}\) and \(\cons{x}{s}\) is harmless, but the
unification of \(\fun{cat}(t,u)\) and~\(t\) is not.

\bigskip

That is why it is best to rename the variables in the pattern so they
are not found in the goal (this is the
\textbf{\(\alpha\)\hyp{}renaming} of \(\lambda\)\hyp{}calculus), and
then try to unify them.

\bigskip

To simplify things even more, we are going to rename \emph{all} the
variables in the definition of `\fun{cat}'.

\end{frame}

% ------------------------------------------------------------------------
%
\begin{frame}{Proving associativity of stack concatenation}

  Let us try for example the following renaming:
  \begin{align*}
    \fun{cat}(        \el,b) &\xrightarrow{\alpha} b\\
    \fun{cat}(\cons{i}{a},b) &\xrightarrow{\smash{\beta}}
    \cons{i}{\fun{cat}(a,b)}
  \end{align*}

  We want to unify the goal \(\fun{cat}(\cons{x}{s},\fun{cat}(t,u))\)
  with the pattern \(\fun{cat}(\cons{i}{a},b)\).

  \bigskip

  The unification yields the equations \(x = i\), \(s = a\) and
  \(\fun{cat}(t,u) = b\).

  \bigskip

  Then we apply the set of equations from the unification, so the
  right\hyp{}hand side of~\(\beta\), \(\cons{i}{\fun{cat}(a,b)}\),
  becomes \(\cons{x}{\fun{cat}(s,\fun{cat}(t,u))}\).

\end{frame}

% ------------------------------------------------------------------------
%
\begin{frame}{Proving associativity of stack concatenation}

  We have:
  \begin{equation*}
    \begin{array}{r@{\;}l@{\;}l@{\qquad}r@{}}
      \ufun{cat}(\cons{x}{s},\fun{cat}(t,u))
      & \xrightarrow{\smash\beta}
      & \cons{x}{\fun{cat}(s,\fun{cat}(t,u))}\\
      & \equiv
      & \cons{x}{\fun{cat}(\fun{cat}(s,t),u)}
      & \IH{\pred{CatAssoc}{s,t,u}}\\
      & \xleftarrow{\smash\beta}
      & \ufun{cat}(\cons{x}{\fun{cat}(s,t)},u)\\
      & \xleftarrow{\smash\beta}
      & \fun{cat}(\ufun{cat}(\cons{x}{s},t),u).
    \end{array}
  \end{equation*}

  Note the use of the inductive hypothesis \(\pred{CatAssoc}{s,t,u}\).

  \bigskip

  This derivation above establishes that
  \(\pred{CatAssoc}{\cons{x}{s},t,u}\) holds.

  \bigskip

  The induction principle entails that \(\forall s,t,u \in
  S.\pred{CatAssoc}{s,t,u}\).\hfill\(\Box\)

\end{frame}

% ------------------------------------------------------------------------
%
\begin{frame}{Unification}

  The best way to think about unification is to think of terms as
  trees, as we did page~\pageref{terms_as_trees}. A function is a
  subtree and a constant is a leaf.

  \bigskip

  The unification algorithm is then very simple:
  \begin{itemize}

    \item if the two terms are constants, they match if they are
      verbatim the same;

    \item if one of the terms is a variable, then they match;

    \item if the two terms are functions, they match if and only if
      \begin{itemize}

        \item they are the same (name),

        \item they have the same \textbf{arity} (number of arguments,
         that is, number of subterms, or number of subtrees),

        \item if they have arguments, the \(n\)-th argument of one
          term matches the \(n\)-th argument of the other.

  \end{itemize}

\end{itemize}

\end{frame}

% ------------------------------------------------------------------------
%
\begin{frame}{Unification / Example \#1}

  For example, let's unify the terms \(\fun{f}(\fun{g}(x,x),
  \fun{h}(y,y))\) and \(\fun{f}(\fun{g}(z), z)\).
  \begin{center}
    \includegraphics{tree-1-1}
  \end{center}
  They do not unify because the subterm \(\fun{g}(x,x)\) does not
  match \(g(z)\): the arity of \fun{g} is different.

  \bigskip

  Note how the terms are non-linear, even though we do not allow them
  in our systems. Indeed, unification is all about equalities, which
  non\hyp{}linearity is about. So let's allow it for exercises.

\end{frame}

% ------------------------------------------------------------------------
%
\begin{frame}{Unification / Example \#2}

  Consider now the terms \(\fun{f}(\fun{g}(x, x), \fun{h}(y, y))\) and
  \(\fun{f}(\fun{g}(\fun{h}(w, \fun{a}()), z), z))\):
  \begin{center}
    \includegraphics{tree-1-2}
  \end{center}

\end{frame}

% ------------------------------------------------------------------------
%
\begin{frame}{Unification / Example \#2}

  The steps are
\begin{align*}
  x &= \fun{h}(w, \fun{a}()) & x &= \fun{h}(w, \fun{a}())
  & \fun{h}(w, \fun{a}()) &= \fun{h}(y, y) & w &= y\\
x &= z      & x &= \fun{h}(y, y) & x      &= \fun{h}(y, y) & x &= \fun{h}(y, y)\\
z &= \fun{h}(y, y) & z &= x      & z      &= x      & y &= \fun{a}()\\
  &                &   &         &        &         & z &= x
\end{align*}
  \noindent Finally we obtain the \textbf{substitution}
\begin{align*}
w &= \fun{a}()\\
x &= \fun{h}(\fun{a}(), \fun{a}())\\
y &= \fun{a}()\\
z &= \fun{h}(\fun{a}(), \fun{a}())
\end{align*}

\end{frame}

% ------------------------------------------------------------------------
%
\begin{frame}{Unification / Example \#2}

  This answer means that if we take the two trees of the query and
  replace the variables according to the substitution, we get the
  exact same tree. Indeed, we get
  \begin{center}
    \includegraphics{tree-1-2u}
  \end{center}

\end{frame}

%% % ------------------------------------------------------------------------
%% %
%% \begin{frame}{Unification / Example \#3}

%%     Consider now the terms \(\fun{f}(\fun{g}(x,x), \fun{h}(\_,\_))\)
%%     and \(\fun{f}(\fun{g}(\fun{h}(w, a()),z), z)\), where the
%%     underscores denote \textbf{fresh variables} (that is, they are
%%     unknown but different from any other variable in the term):
%%     \begin{center}
%%       \includegraphics{tree-1-3}
%%     \end{center}

%% \end{frame}

%% % ------------------------------------------------------------------------
%% %
%% \begin{frame}{Unification / Example \#3}

%%   The steps are now
%%   \begin{align*}
%%     x &= \fun{h}(w,\fun{a}()) & x &= \fun{h}(w, \fun{a}())
%%       & x &= \fun{h}(w, \fun{a}()) & x &= \fun{h}(\alpha, \fun{a}())\\
%%     x &= z & x &= z & z &= x & z &= \fun{h}(\alpha, \fun{a}())\\
%%     \fun{h}(\alpha, \beta) &= z &
%%     \fun{h}(\alpha, \beta) &= \fun{h}(w, \fun{a}())
%%     & w &= \alpha & w &= \alpha\\
%%     & & & & & & \beta &= \fun{a}()
%%   \end{align*}

%%   Finally
%%   \begin{align*}
%%     \beta &= \alpha \\
%%     w &= \alpha\\
%%     x &= \fun{h}(\alpha, \fun{a}())\\
%%     z &= \fun{h}(\alpha, \fun{a}())
%%   \end{align*}

%% \end{frame}

%% % ------------------------------------------------------------------------
%% %
%% \begin{frame}{Unification / Example \#3}

%%   This means that if we replace the variables in the two initial
%%   trees, we get the same tree. Indeed, we get
%%   \begin{center}
%%     \includegraphics{tree-1-3u}
%%   \end{center}

%% \end{frame}

% ------------------------------------------------------------------------
%
\begin{frame}{More proofs by induction}

  Here is again the concatenation of two stacks:
  \begin{align*}
    \fun{cat}(        \el,t) &\xrightarrow{\alpha} t\\
    \fun{cat}(\cons{x}{s},t) &\xrightarrow{\smash{\beta}}
    \cons{x}{\fun{cat}(s,t)}
  \end{align*}
  Here is a functional rewrite system that defines a function
  `\fun{len}' that computes the length of a stack.
  \begin{equation*}
    \begin{array}{r@{\;}l@{\;}l@{}}
      \fun{len}\,\el & \smashedrightarrow{\gamma} & 0;\\
      \fun{len}(\cons{x}{s}) & \smashedrightarrow{\delta}
      & 1 + \fun{len}\, s.
    \end{array}
  \end{equation*}
  So, for example: \(\fun{len}\,[1;2;3] \twoheadrightarrow 1 + (1 + (1
  + (0))) = 3.\)

  \bigskip

  Let us prove the property
  \begin{equation*}
    \pred{LenCat}{s,t} \colon \fun{len}(\fun{cat}(s, t)) \equiv
    \fun{len}\,s + \fun{len}\,t.
  \end{equation*}
  That is: concatenation preserves the length of the stacks.

\end{frame}

% ------------------------------------------------------------------------
%
\begin{frame}{More proofs by induction}

  Note how there are \textbf{two parameters}, \(s\)~and~\(t\), instead
  of one.

  \bigskip

  In this case, we can try to prove the property by induction on one
  parameter, as we did before, and hope that the induction hypothesis
  with the other parameter considered constant suffice. Let us chose
  the first one.

  \bigskip

  The base case is \(\forall t \in S.\pred{LenCat}{\el,t}\), that is
  \[\fun{len}(\fun{cat}(\el,t)) \equiv \fun{len}\,\el +
  \fun{len}\,t.\] It is immediate (the call being rewritten is
  underlined):
  \begin{equation*}
    \fun{len}(\ufun{cat}(\el,t)) \smashedrightarrow{\alpha}
    \fun{len}\,t = 0 + \fun{len}\,t
    \smashedleftarrow{\gamma} \ufun{len}\,\el + \fun{len}\,t.
  \end{equation*}

\end{frame}

% ------------------------------------------------------------------------
%
\begin{frame}{More proofs by induction}

  The general case \(\forall s,t \in S.\pred{LenCat}{s,t} \Rightarrow
  \forall x \in T.\pred{LenCat}{\cons{x}{s}, t}\), where \(S\)~is the
  set of all stacks, and \(T\)~is the set of all terms (\(S \subset
  T\)).

  \bigskip

  Let us assume now \(\forall s,t \in S.\pred{LenCat}{s,t}\) (the
  induction hypothesis), and prove \(\forall x \in
  T.\pred{LenCat}{\cons{x}{s}, t}\).

  \bigskip

  We have:
  \begin{equation*}
    \begin{array}{r@{\;}l@{\;}l@{\qquad}r@{}}
        \fun{len}(\ufun{cat}(\cons{x}{s},t))
      & \smashedrightarrow{\beta}
      & \ufun{len}(\cons{x}{\fun{cat}(s,t)})\\
      & \smashedrightarrow{\delta}
      & 1 + \fun{len}(\fun{cat}(s,t))\\
      & \equiv & 1 + (\fun{len}\,s + \fun{len}\,t)
      & \pred{LenCat}{s,t}\\
      & = & (1 + \fun{len}\,s) + \fun{len}\,t
      & \text{associativity of \((+)\)}\\
      & \smashedleftarrow{\delta}
      & \ufun{len}(\cons{x}{s}) + \fun{len}\,t.
    \end{array}
  \end{equation*}
  We proved \(\forall x \in T.\pred{LenCat}{\cons{x}{s}, t}\).\hfill\(\Box\)

\end{frame}

% ------------------------------------------------------------------------
%
\begin{frame}{Higher-order rewrite systems}

  We already extended term-rewriting systems (TRS) with conditions.

  \bigskip

  We can also extend them in another direction: variables could denote
  functions. These are called \textbf{higher\hyp{}order systems}.

  \bigskip

  They enable defining functions that can be parameterised by other
  functions.

  \bigskip

  For example, let us consider \textbf{iterators}. Typically, those
  are functions that traverse a data structure and apply a given
  function at each step.

\end{frame}

% ------------------------------------------------------------------------
%
\begin{frame}{Maps on stacks}

  In the case of stacks, a simple iterator is a function `\fun{map}'
  that takes a stack and returns a stack made of the images of the
  elements, in the original order. So $[e_1; e_2; \ldots; e_n]$ becomes $[f(e_1);
    f(e_2); \ldots; f(e_n)]$.
  \begin{equation*}
    \begin{array}{r@{\;}l@{\;}l@{}}
      \fun{map}(f, \el) &\smashedrightarrow{\alpha} & \el;\\
      \fun{map}(f, \cons{x}{s}) &\smashedrightarrow{\beta}
      & \cons{f(x)}{\fun{map}(f, s)}.
    \end{array}
  \end{equation*}

  For example, if we set
  \begin{align*}
    \fun{incr}(n)   &\rightarrow n + 1.\\
    \fun{square}(n) &\rightarrow n * n.
  \end{align*}
  then we have
  \begin{align*}
    \fun{map}(\fun{incr}, [4; 4; 5; 0]) &\twoheadrightarrow [5; 5; 6; 1].\\
    \fun{map}(\fun{square}, [4; 4; 5; 0]) &\twoheadrightarrow [16; 16; 25; 0].
  \end{align*}

\end{frame}

% ------------------------------------------------------------------------
%
\begin{frame}{Folds on stacks}

  Folds take an additional parameter with is supposed to be the
  initial value of an \textbf{accumulator}. In this framework, the
  given function is applied to an element of the stack \emph{and to the
  current value of the accumulator}.

  \bigskip

  For example, for summing all integers in a stack, we would have
  \begin{equation*}
    \fun{fold} (\fun{add}, [4; 4; 5; 0], 0) \twoheadrightarrow 13,
  \end{equation*}
  where $\fun{add}(x,y) \rightarrow x + y$.

  \bigskip

  Note how the direction of iteration does not matter, because
  \begin{equation*}
    4 + (4 + (5 + (0 + \underline{0})))
    = (((\underline{0} + 4) + 4) + 5) + 0,
  \end{equation*}
  where the initial value of the accumulator is~\(\underline{0}\).

\end{frame}

% ------------------------------------------------------------------------
%
\begin{frame}{Lambdas}

  Note how cumbersome it is sometimes to have to define the function
  applied by the iterator (here a fold), when it is probably used only
  once.

  \bigskip

  To remedy this, we can further extend or systems with
  \textbf{anonymous functions}, also known as \textbf{lambdas} (a
  name coming from the $\lambda$\hyp{}calculus, which is a different
  framework in which to model computations.)

  \bigskip

  For example, we will write $\Xfun\; x \rightarrow x$ to denote the
  identity function, without naming it. Of course, we have to forbid
  calling any function `\fun{fun}` now (in other words, \Xfun is a
  \textbf{keyword}).

  \bigskip

  In our case above:
  \begin{equation*}
    \fun{fold} ((\Xfun\; (x,y) \rightarrow x + y), [4; 4; 5; 0], 0)
    \twoheadrightarrow 13,
  \end{equation*}

\end{frame}

% ------------------------------------------------------------------------
%
\begin{frame}{A map is a fold}

  A map is then a special case of a fold, where the accumulator is the
  new stack:
  \begin{equation*}
    \fun{map}(f, s)
    \rightarrow \fun{fold}((\Xfun\; (e, a) \rightarrow \cons{f(e)}{a}), s, \el),
  \end{equation*}
  where $a$ is the accumulator and $e$ a stack element.

  \bigskip

  Contrary to $(+)$, pushing on a stack is not a symmetric relation,
  so the order of traversal matters.  If `\fun{fold}' visited the
  stack elements from left to right, then the result of the call
  $\fun{map}(\fun{incr}, [4; 4; 5; 0])$ would be $[1;6;5;5]$, which is
  the reverse of what we expect, which is $[5;5;6;1]$.

  \bigskip

  Here, since we want to conserve the order of the original elements
  of the stack, the fold must apply the given function first to the
  bottom element, then the penultimate, then the antepenultimate
  etc. until the top. If we consider a stack a being layout
  horizontally, this means applying the function from right to left.

  \bigskip

  This is called a \textbf{right fold}.

\end{frame}

% ------------------------------------------------------------------------
%
\begin{frame}{Right fold}

  Let's rename `\fun{fold}' as `\fun{foldr}' (`fold right'). Then:
  \begin{equation*}
    \begin{array}{r@{\;}l@{\;}l@{}}
      \fun{foldr} (f, [], a) & \rightarrow a;\\
      \fun{foldr} (f, \cons{x}{s}, a) & \rightarrow f(x, \fun{foldr}(f, s, a)).
    \end{array}
  \end{equation*}
  A way to understand this definition is by extension:
  \begin{equation*}
    \fun{foldr}(f,[e_1; e_2; \ldots; e_n], a) \equiv f(e_1, f(e_2, \ldots f(e_n, a))).
  \end{equation*}

  Note that the recursive call in the definition of `\fun{foldr}' is
  not in tail position. (See page~\pageref{tail_call} an onwards.)

  \bigskip

  In this case, this means that all the calls to the iterated
  function~$f$ will be waiting for all the calls to~`\fun{foldr}' to
  be finished (as well as all the calls to~$f$ on the previous
  elements of the stack).

  \bigskip

  Generally speaking, the call stack will grow proportionally to the
  length of the stack.

\end{frame}

% ------------------------------------------------------------------------
%
\begin{frame}{Left fold}

  The dual iteration of a right fold is a \textbf{left fold}:
  \begin{equation*}
    \fun{foldr}(f,[e_1; e_2; \ldots; e_n], a) \equiv f(e_n, \ldots f(e_2, f(e_1, a))).
  \end{equation*}
  It is a simple matter of rearranging the calls in the definition of `\fun{foldr}`:
  \begin{equation*}
    \begin{array}{r@{\;}l@{\;}l@{}}
      \fun{foldl} (f, [], a) & \rightarrow a;\\
      \fun{foldl} (f, \cons{x}{s}, a) & \rightarrow \fun{foldl}(f, s, f(x,a)).
    \end{array}
  \end{equation*}

  Note how the recursive call is in tail position, so, if tail call
  optimisation is performed, the call stack will grow a constant a
  constant amount if the calls to the given function also grow the
  stack a constant amount.

  \bigskip

  These considerations on memory usage are only meaningful when
  running our programs, but not when designing a system.

\end{frame}

% ------------------------------------------------------------------------
%
\begin{frame}{Binary trees}

  \begin{itemize}

    \item A \textbf{binary tree} is a special case of tree, where each
      node has either no children or exactly two.
      \begin{center}
        \includegraphics{bt_shape}
      \end{center}

    \item Nodes are of two kinds: \textbf{internal}
      ({\tiny\raisebox{1pt}{\(\oslash\)}}, \(\circ\) and \(\bullet\))
      or \textbf{external} (\(\scriptstyle \Box\)).

    \item The \textbf{root} is the only internal node without a parent
      ({\tiny\raisebox{1pt}{\(\oslash\)}}).

    \item \textbf{Leaves} (\(\bullet\)) are internal nodes whose
      children are two external nodes (\(\scriptstyle \Box\)).

  \end{itemize}

\end{frame}

% ------------------------------------------------------------------------
%
\begin{frame}{Binary trees}

  Internal nodes are usually associated with some kind of information,
  whilst external nodes are not, like
  \begin{center}
    \includegraphics[bb=71 645 190 721]{bt_ex1}
    \qquad
    \includegraphics[bb=71 645 190 721]{bt_ex2}
  \end{center}
  \centerline{Extended binary tree \qquad\quad Pruned binary tree}

  We need to ponder how many shapes a binary tree can take, just like
  stacks. Each shape as a value constructor and maybe some data.
  \begin{itemize}

    \item The \textbf{empty tree} is an external node, without
      information: \(\fun{ext}().\)

    \item A \textbf{non\hyp{}empty tree} is an internal node carrying
      some information \(x\), and which has a left child~\(t_1\) and a
      right child~\(t_2\): \(\fun{int}(x, t_1, t_2).\)

  \end{itemize}

\end{frame}
